\setchapterpreamble[u]{\margintoc}
\chapter{T-SQL: Visões}
\labch{cet615-visoes}

Este capítulo cobre a apresentação sobre visões da segunda
aula\sidecite{BaiocoCET615}. No SSMS em português, as visões aparecem
como \emph{Exibições}.

\section{O que é uma visão}

Uma \textbf{visão}\index{visão} (\emph{view}) é uma consulta guardada
com nome, que pode ser usada como se fosse uma tabela. A imagem da aula
é a de uma moldura: através dela se veem só os dados de interesse de uma
ou mais tabelas (ou de outras visões). O SGBD guarda apenas a definição;
os dados continuam nas tabelas. Uma mudança feita nas tabelas aparece
na visão, e vice-versa.

\subsection{Vantagens}

\begin{itemize}
	\item \textbf{Segurança.} A visão pode esconder colunas e linhas. Em
	aula, a professora foi além do slide: em vez de dar ao usuário
	permissão nas três tabelas de uma consulta, dá-se permissão só na
	visão. O acesso passa a ser por funcionalidade, e o usuário nunca toca
	as tabelas.
	\item \textbf{Simplicidade.} Uma consulta complexa e frequente vira um
	\texttt{select * from visao}.
	\item \textbf{Valores calculados e resumos} ficam prontos na visão.
\end{itemize}

\subsection{Desvantagens}

\begin{itemize}
	\item O desempenho é o do \texttt{select} que define a visão. Se a
	visão tem junção, a junção roda mesmo quando a consulta só precisa de
	colunas de uma das tabelas.
	\item Alterar dados por meio de visões tem restrições (veja adiante).
\end{itemize}

A professora corrigiu a leitura do slide: a visão em si não é ruim para
o desempenho. O problema é usá-la quando a junção dela não é necessária.

\section{Criar, usar, alterar e excluir}

\begin{lstlisting}[caption={Visão sobre o banco de biblioteca.}]
create view emprestimos_detalhados (leitor, livro, data_emprestimo)
as
select le.nome, l.titulo, e.data
from leitor le inner join emprestimo e
     on le.codleitor = e.codleitor
     inner join livro l
     on e.codlivro = l.codlivro
go

select * from emprestimos_detalhados
where data_emprestimo >= '2026-11-01'
\end{lstlisting}

\begin{itemize}
	\item A visão é criada no banco em uso, e o \texttt{create view} precisa
	ser o primeiro comando do lote (\refch{cet615-sqlserver}).
	\item Os nomes das colunas vão na lista do \texttt{create view} ou como
	apelidos no \texttt{select}. São obrigatórios quando duas colunas
	teriam o mesmo nome (dois \texttt{nome} numa junção) e em colunas
	calculadas.
	\item Para criar uma visão, basta saber escrever o \texttt{select}. A
	recomendação é testar o \texttt{select} antes de transformá-lo em
	visão.
	\item O uso ideal é \texttt{select * from visao}, com \texttt{where}
	quando preciso. Um critério que muda com o tempo, como um intervalo de
	datas, fica fora da visão, no \texttt{where} de quem a consulta.
\end{itemize}

\begin{lstlisting}
alter view emprestimos_detalhados (leitor, livro, data_emprestimo)
as
select ...   -- o select inteiro precisa ser reescrito

exec sp_helptext emprestimos_detalhados   -- mostra a definição

drop view emprestimos_detalhados
\end{lstlisting}

O \texttt{alter view} substitui a definição e mantém as permissões já
concedidas. O \texttt{drop} seguido de \texttt{create} perde as
permissões. Por isso o \texttt{alter} só faz sentido com o sistema em
produção; em projeto, apagar e recriar resolve.

\subsection{O antipadrão demonstrado em aula}

A professora mostrou de propósito o uso errado. Para listar só os nomes
dos fornecedores, ela consultou a visão de fornecedores e produtos. O
resultado veio com nomes repetidos, sinal de que a junção estava sendo
executada. Um \texttt{distinct} ``resolve'', mas a forma certa é
consultar a tabela \texttt{fornecedor} diretamente. Já o filtro pelo
nome do produto faz sentido na visão, porque a junção é necessária para
chegar a ele.

\begin{lstlisting}
-- antipadrão: junção desnecessária, remendada com distinct
select distinct leitor from emprestimos_detalhados

-- certo
select nome from leitor
\end{lstlisting}

\section{Alterar dados por meio de visões}

Segundo os slides, uma visão só é atualizável se o seu \texttt{select}:

\begin{itemize}
	\item não tem funções de agregação na lista de seleção;
	\item não tem colunas calculadas a partir de outras;
	\item referencia pelo menos uma tabela no \texttt{from}. Uma visão como
	\texttt{select getdate() as hoje} nunca é atualizável.
\end{itemize}

E os comandos têm restrições:

\begin{itemize}
	\item \texttt{insert} e \texttt{update} só podem alterar uma das
	tabelas do \texttt{from};
	\item o \texttt{insert} precisa satisfazer todas as colunas
	\texttt{not null} das tabelas, mesmo as que não aparecem na visão.
	Como a visão é só uma moldura, a coluna obrigatória que ficou de fora
	fica sem valor, e a inserção falha;
	\item \texttt{delete} não é permitido em visão com mais de uma tabela.
\end{itemize}

A conclusão da professora é que \emph{visão serve para consulta}. Para
usar \texttt{insert} ou \texttt{update} por uma visão, o usuário
precisaria conhecer o \texttt{select} dela para saber em que tabela está
mexendo. E incluir na visão todas as colunas obrigatórias, para o
\texttt{insert} funcionar, desfaz o objetivo de segurança. Para escrita,
há saídas melhores: outro usuário com permissão nas tabelas ou um
procedimento armazenado (\refch{cet615-procedimentos}).

\subsection*{Complemento: como o SGBD executa uma visão}

\textcite[seção 5.3.3]{Elmasri2011} descreve duas estratégias de
implementação. Na \textbf{modificação de consulta}, o SGBD reescreve a
consulta sobre a visão como uma consulta sobre as tabelas de base. Na
\textbf{materialização}, ele grava o resultado da visão numa tabela e a
mantém atualizada de forma incremental enquanto ela é usada.

No SQL Server, a visão comum segue a primeira estratégia: a
documentação\sidecite{MSCreateView} diz que uma operação sobre a visão é
validada e traduzida em uma operação sobre as tabelas de base. Só a
\emph{visão indexada} guarda dados. Isso reforça o ponto da professora:
o custo de usar uma visão é o custo do \texttt{select} que ela esconde.

O livro também explica por que a atualização por visão é restrita: numa
visão com junção, uma mesma alteração pode corresponder a várias
alterações diferentes nas tabelas de base, e o SGBD não tem como saber
qual foi a pretendida.

\subsection*{Complemento: opções do \texttt{create view}}

A documentação\sidecite{MSCreateView} traz recursos que a aula não
cobre:

\begin{description}
	\item[\texttt{with check option}:] impede que uma alteração feita pela
	visão produza uma linha que a própria visão deixaria de mostrar.
	\item[\texttt{with schemabinding}:] amarra a visão às tabelas; elas não
	podem ser alteradas de um jeito que afete a visão, nem apagadas antes
	dela. É pré-requisito da visão indexada.
	\item[\texttt{order by}:] não é aceito na definição, a não ser com
	\texttt{top}, e mesmo assim não garante a ordem do resultado. A ordem
	vai na consulta sobre a visão.
	\item[Gatilhos \texttt{instead of}:] tornam atualizável uma visão que
	não seria, porque o gatilho decide o que fazer com cada alteração
	(\refch{cet615-gatilhos}).
\end{description}

A documentação também atualiza um detalhe do slide: o texto da visão,
que o \texttt{sp\_helptext} mostra, fica na visão de catálogo
\texttt{sys.sql\_modules}; \texttt{syscomments} é uma visão de
compatibilidade.
